%% content/faq.tex -- extra questions and answers for the website's FAQ page (parents/faq/).
%%
%% The FAQ page shows the questions from the presenter notes in main.tex (the third argument of
%% \kcnotes) first, then the ones below. This file is NOT part of the deck, the presenter notes or
%% the handout; it holds the questions families ask outside the information evening, so the slides
%% stay short. tools/site/parents.py reads it.
%%
%% Format:
%%   \faqtopic{Name}      starts (or adds to) a topic; use a name from TOPIC_ORDER in
%%                        site/src/pages/parents/faq.astro, or a new one (it is listed last)
%%   \qa{Question}{Answer}
%% Macros from content/shared.tex work (\LevelWeeks, \Homework, ...). \textcite{key} and
%% \parencite{key} link to the reference list (keys from references.bib). \sitelink{path/}{text}
%% links to a page on this website (path without a leading slash).
%%
%% Course facts must agree with docs/SYLLABUS.md, docs/DESIGN.md and docs/ASSUMPTIONS.md.
%% Research claims must use only what NOTES.md lists for that source. Written November 2026.

%% =====================================================================
\faqtopic{Before you join}

\qa{Who is the course for?}{Children in Grades 4 and 5 (about 9 to 11 years old) who are curious
  about how computers work. No experience is needed. Everything children read is written at a
  Grade 4--5 reading level, with pictures, worked examples and a glossary.}

\qa{Isn't 9 or 10 too young to learn to code?}{No. Scratch, the tool we start with, was designed at
  MIT for children of this age \parencite{resnick2009scratch}. Snap-together blocks mean there are no
  spelling mistakes to fight, so children can focus on the ideas. We move to typed Python only in
  week 7 of Level 1, once those ideas feel familiar.}

\qa{My child has never coded. Is that OK?}{Yes. Level 1 assumes nothing at all and there is no
  entry test: instead, children fill in a short, ungraded survey about what they enjoy and what
  feels new. The first lessons are about giving clear step-by-step instructions, which every child
  already does every day.}

\qa{My child already knows some Scratch. Should they skip Level 1?}{Usually not. Level 1 moves on
  to Python halfway through, and every lesson has challenge tasks for children who finish early.
  If your child has already written Python programs, talk to the course leader: Levels 2 and 3
  start with a short placement quiz that shows where to begin.}

\qa{What math does my child need?}{Only what Grade 4--5 children already meet at school: whole
  numbers, division with remainders, simple fractions and decimals, angles in degrees and
  coordinates on a grid. When an older idea comes up (for example negative numbers), the workbook
  explains it in a short \enquote{Math helper} box first.}

\qa{How do I know if my child is ready?}{Interest matters more than skill. A child who enjoys
  building, puzzles, drawing or making up games, and can stay with a task for about 15 minutes, is
  ready. Reading help from an adult is fine; pairs also read instructions to each other.}

\qa{Can we look at the course before we sign up?}{Yes. Every workbook chapter, homework set and
  starter file is free on this website. Try the week 1 chapter of
  \sitelink{levels/level-1/}{Level 1} together and see how your child responds.}

\qa{How much does it cost?}{The curriculum, the website and all the apps (Scratch, Thonny and
  Python) are free and need no subscription. If there is a fee for the sessions themselves, the
  person who runs your group will tell you.}

%% =====================================================================
\faqtopic{The course}

\qa{How long is the course?}{Three levels of \LevelWeeks{} each, with \SessionLength{}, so about
  \TotalWeeks{} in all. Holiday weeks without a session do not count. Each level ends with a
  finished project, so your child can stop after any level.}

\qa{What happens in a session?}{Every session has the same five parts, so children always know
  what comes next: a warm-up game, a short new idea, a guided demonstration, a hands-on build in
  pairs, and a wrap-up where children share and reflect. The group is \GroupSize{}, with a screen
  break in every session.}

\qa{What if my child misses a session?}{Tell the course leader. Each week's workbook chapter,
  homework and starter files are on this website, so your child can catch up at home in about
  half an hour. Week 6 of every level is a review week, and each level has spare sessions that
  can be used for catching up.}

\qa{Are there tests or grades?}{No grades. Each level has a short mid-level checkpoint and an
  end-of-level check (Levels 2 and 3 also start with a placement quiz), but they only help the
  adult see what to revisit. Projects are scored on a friendly four-step scale, from \enquote{Just
  starting} to \enquote{Expert}, which values creativity, working code and being able to explain
  it equally.}

\qa{What happens at the end of each level?}{Children spend two sessions on their own final project
  and then show it at a showcase for families in the last week. Everyone gets a certificate,
  filled in by hand. The showcase happens on screen in the room; nothing is uploaded.}

\qa{Does the course follow any standards?}{Yes. Lessons are mapped to the CSTA K--12 Computer
  Science Standards \parencite{csta2017standards} and to the five big ideas of the AI4K12
  guidelines \parencite{touretzky2019ai4k12}. Level 3 borrows a few high-school standards on
  purpose, for children ready to stretch.}

%% =====================================================================
\faqtopic{What your child learns}

\qa{What will my child be able to make?}{Level 1 ends with \emph{\LevelOneProject}:
  \LevelOneProjectLine. Level 2 ends with an \emph{\LevelTwoProject}: \LevelTwoProjectLine.
  Level 3 ends with \emph{\LevelThreeProject}: \LevelThreeProjectLine.}

\qa{Why start with Scratch and then switch to Python?}{Blocks let children learn the big ideas
  (steps, loops, choices, variables) without spelling mistakes getting in the way. Studies find
  that children who start with blocks are ahead at first and do just as well once they move to
  typed code \parencite{weintrop2019transition}. Python comes next because it reads almost like
  English and is used in schools, universities and real jobs.}

\qa{Why Python and not JavaScript, Java or the language of a favorite game?}{Python has very little
  punctuation to get wrong, and Thonny, the editor we use, was made for beginners and shows
  clearly what the program is doing. It works fully offline, with no accounts. The ideas children
  learn (loops, functions, lists, objects) carry over to every other language.}

\qa{What is \enquote{debugging}, and why does the course talk about it so much?}{Debugging means
  finding and fixing mistakes in a program. Every session includes a deliberate bug to hunt,
  because all programmers make mistakes and the skill is finding them calmly. Teaching a clear
  debugging routine has been shown to improve both debugging skill and children's confidence
  \parencite{michaeli2019debugging}.}

\qa{Does the course teach artificial intelligence?}{Yes, in Level 3, and from the inside: children
  build tiny AI programs by hand in Python (for example one that guesses by finding the most
  similar example, and one that makes up silly sentences from patterns), so they see that AI
  learns from data and can be wrong. No online AI services are used.}

\qa{Will my child also learn about online safety?}{Yes. Every level has at least five short
  conversations about being a good digital citizen: keeping personal information private, giving
  credit for other people's work, making things everyone can use, asking for consent and using AI
  fairly.}

\qa{My child prefers to work alone. Do they have to work in a pair?}{For the main build, yes,
  because children explain their thinking out loud in pairs, and studies find pairs often learn
  more, especially less-experienced children \parencite{denner2014pair}. Roles switch often, and
  there is also free-build time to work on their own ideas. Give it a few weeks; many children
  who are unsure at first come to enjoy it.}

%% =====================================================================
\faqtopic{Computers and apps}

\qa{What software does my child use, and does it cost anything?}{Three free tools: the offline
  Scratch app from MIT, Thonny (a beginner's Python editor with Python built in) and, in Level 3,
  git, which Apple provides for free. None of them needs an account or a subscription.}

\qa{How do I set up our Mac?}{Follow the \sitelink{resources/setup-mac/}{Mac setup guide}. Allow
  about an hour, once, before the first session. You can copy one command that installs
  everything, or install each app by hand step by step. The internet is needed only for this
  one-time installation, and a grown-up should do it.}

\qa{We have a Windows PC or a Chromebook. Can my child still practice at home?}{The course is
  written and tested for Macs, so menu names and keyboard shortcuts in the materials are the Mac
  ones. The Scratch app is also free on Windows and Chromebooks, and Thonny is free on Windows, so
  most homework works there with small differences. The paper parts of the homework work
  anywhere.}

\qa{Can my child use an iPad or tablet?}{Not for this course. There is no Scratch app for iPad,
  and Thonny needs a computer with a keyboard. A tablet is fine for reading the workbook on this
  website.}

\qa{Does my child need the internet at home?}{No. Lessons and homework work completely offline.
  The internet is needed only once, to install the apps.}

\qa{Does my child need their own computer?}{No. In sessions, pairs share one Mac. At home, the
  homework (\Homework) can be done on a family computer.}

\qa{Something on our Mac is not working. Where can we get help?}{Start with the
  \sitelink{resources/troubleshooting/}{troubleshooting guide}, which covers the most common
  problems with Scratch and Thonny. In Thonny, the menu \emph{View \textgreater{} Assistant}
  explains error messages in friendlier words. If you are still stuck, ask the course leader.}

\qa{Where are the answers? Can I check my child's work?}{Each Python exercise on this website has
  three hints that get more and more specific, and many exercises come with a self-check program
  that says PASS or TRY AGAIN with a tip. Full answers stay with the course leader, so children
  can enjoy solving the problems themselves.}

\qa{Can my child keep the projects they make?}{Yes. Projects are saved as files on the Mac. You can copy them to a USB stick to keep. In Level 3 children also learn git,
  which keeps a full history of every version of their project.}

%% =====================================================================
\faqtopic{Helping at home}

\qa{I can't code. Can I still help?}{Yes, and you do not need to learn. The most useful things are
  asking your child to show you what they made, asking about the bugs they found, and praising the
  way they worked rather than \enquote{being a natural}. Being curious and nearby, without taking
  over the keyboard, helps more than knowing the answer.}

\qa{How much homework is there, and what if we skip a week?}{About \Homework. There is a short Core
  part for everyone, plus optional Challenge and Creative parts for children who want more.
  Skipping a week now and then is fine; a regular, protected time each week helps most.}

\qa{How much should I help with homework?}{Ask questions instead of giving answers: \enquote{What
  did you expect to happen? What happened instead? What could you try?} A little struggle is part
  of learning. If your child is stuck for more than about ten minutes, suggest a hint from the
  website, or write down the question for the next session.}

\qa{My child gets upset when the code doesn't work. What should I do?}{It is completely normal;
  every programmer meets bugs every day. Suggest a strategy: read the error message, read the code
  out loud, change one thing at a time, or take a short break. Praise the strategy, not just the
  fix, and try the word \enquote{yet}: \enquote{It doesn't work \emph{yet}.}}

\qa{What should I ask my child after a session?}{Good questions are \enquote{Show me what you
  made}, \enquote{What was the trickiest bug?}, \enquote{How did you find it?} and \enquote{What
  would you add next?} Each level also has a short note for parents with ideas for that level.}

\qa{My child uses words I don't understand. Where can I look them up?}{Each level has a glossary
  and a one-page cheat sheet on this website, written for children, so they are quick for adults
  to read too. Asking your child to explain a word is even better: explaining is one of the best
  ways to learn.}

\qa{My child wants to code more than once a week. What can they do?}{Wonderful. The Challenge and
  Creative parts of each homework, and the extra exercises on this website, give plenty to do. The
  Scratch app is also great for free play: suggest making a small game or story of their own.
  Keep it within your family's screen-time plan.}

\qa{Should I learn to code alongside my child?}{Only if you want to, but many families enjoy it. The
  workbook is short and friendly for adults too. Letting your child be the teacher, while you ask
  questions, is a great way to share it.}

%% =====================================================================
\faqtopic{Screen time and wellbeing}

\qa{Isn't this just more screen time?}{It is a small amount: \SessionLength{} plus about 15 minutes
  of homework. Every session has a screen break and an activity away from the computer, and the
  screen time is creative, not passive. Pediatricians recommend that families make their own media
  plan rather than follow one fixed number of hours \parencite{aap2016media}; this course can fit
  into that plan.}

\qa{My child already plays a lot of video games. Will this make it worse?}{Making a game is very
  different from playing one: it means planning, testing and fixing. Many children start to see
  the games they play as something people built, which they can understand and question. It is
  still screen time, so count it in your family plan.}

\qa{How does the course look after children's eyes and bodies?}{Sessions are one hour with a break,
  children stretch and look away from the screen regularly (the 20-20-20 rule: every 20 minutes,
  look at something 20 feet away for 20 seconds), and some activities happen with cards, paper or
  moving around.}

%% =====================================================================
\faqtopic{Safety and privacy}

\qa{Will my child need an email address or online account?}{No. No child creates an account,
  logs in or gives an email address, at any level. US law requires parental consent before online
  services collect personal information from children under 13 \parencite{ftc2025coppa}; this
  course simply never collects any.}

\qa{I've heard about the Scratch online community. Is it safe?}{We use the offline Scratch app, so
  children never join or see the online community, and nothing they make is shared. If your child
  later wants to share projects online, read Scratch's guide for parents together first and decide
  as a family.}

\qa{What information about my child does the course collect?}{None. Programs ask only for
  nicknames or favorite colors, git in Level 3 uses a nickname, the example data is made up, and
  projects stay on the computers in the room. This website has no accounts, no analytics and no
  cookies.}

\qa{Does the course use ChatGPT or other AI chat tools?}{No. The course does not use any online AI
  services, because they need accounts and send what children type to other companies. In Level 3
  children build tiny AI programs themselves, which teaches more about how AI works than using a
  chatbot.}

%% =====================================================================
\faqtopic{Every learner}

\qa{My child has dyslexia or finds reading hard. Can they take part?}{Yes. The student
  workbooks and homework come in a dyslexia-friendly edition with a clearer font, bigger text, more spacing and a cream
  page, and this website has an \emph{Easy-read font} switch (see
  \sitelink{about/\#accessibility}{accessibility}). Instructions are short, Scratch needs very
  little typing, and pairs read aloud to each other.}

\qa{My child has ADHD or finds it hard to sit still. Will this work?}{Many children with ADHD love
  coding because they see results quickly. The course helps with a predictable routine (the same
  five parts each session), short chunks, a screen break, and activities that involve moving
  around. Tell the course leader what works for your child at school or at home.}

\qa{My child is color-blind. Will that be a problem?}{No. All materials use a color palette that is
  safe for color-blind readers \parencite{okabe2008cud}, and every colored box also has an icon and
  a word, so no meaning depends on color alone.}

\qa{My child is very advanced. Will they be bored?}{Every lesson has an extension route, and
  homework has optional Challenge and Creative parts. Final projects are open-ended, so a fast
  learner can go much further. Level 3 covers ideas often met in high school, such as recursion,
  sorting and how simple AI learns.}

\qa{English is not our first language at home. Can my child still take part?}{Yes. Materials use
  short sentences at a Grade 4--5 reading level, with a glossary for each level. The few English
  words in Python (such as \texttt{print} or \texttt{if}) are learned like vocabulary. You can help
  in any language by asking your child to explain what their program does.}

%% =====================================================================
\faqtopic{Why learn to code now}

\qa{If AI can write code, why should my child learn to code?}{Someone still has to understand, check
  and fix what AI writes, and decide what to build in the first place. Research is still early:
  in one study of 69 learners aged 10 to 17, an AI helper let them finish more tasks without making
  them worse at editing code by hand \parencite{kazemitabaar2023codex}, while a 2026 OECD report
  warns that general chatbots can raise scores without lasting learning \parencite{oecd2026deo}.
  Knowing how to code helps children use AI as a tool rather than depend on it.}

\qa{What is computational thinking?}{It is a way of solving problems: breaking a big problem into
  smaller parts, spotting patterns, ignoring details that do not matter, and writing clear steps
  \parencite{wing2006ct}. Coding is one of the best ways to practice it, and the same habits
  help with planning projects in everyday life.}

%% =====================================================================
\faqtopic{After the course}

\qa{What can my child do after Level 3?}{Keep building! Children leave Level 3 able to plan, write,
  test and save larger programs on their own. Good next steps are computer science classes or
  clubs at school, longer projects of their own in Python, or free project guides from
  organizations such as the Raspberry Pi Foundation or Code.org. If a site needs an account,
  read its privacy rules and set it up together.}

\qa{Will this help with computer science at school later?}{It gives a strong start: the ideas in
  this course (loops, functions, lists, objects, testing) are the same ones met in middle- and
  high-school computer science, and Python is widely used there. It is not a shortcut past those
  classes, but children usually arrive confident and curious.}
